\begin{tikzpicture}[
  font=\small,
  input/.style={draw, rounded corners=2pt, align=center, minimum height=9mm,
                text width=13mm, fill=gray!7, inner sep=2pt},
  op/.style={draw, rounded corners=2pt, align=center, minimum height=10mm,
             text width=23mm, inner sep=2pt},
  merge/.style={draw, rounded corners=2pt, align=center, minimum height=12mm,
                text width=21mm, inner sep=2pt},
  read/.style={draw, rounded corners=2pt, align=center, minimum height=10mm,
               text width=17mm, fill=gray!8, inner sep=2pt},
  arrow/.style={-{Latex[length=2.1mm]}, thick}
]
% Panel (a): GraphConv additive fusion.
\node[font=\bfseries,anchor=west] at (-6.20,1.45) {(a) GraphConv 加法融合基线};
\node[input] (xi) at (-5.60,0.55) {$\mathbf x_i$};
\node[input] (xj) at (-5.60,-0.55) {$\{\mathbf x_j\}$};
\node[op,fill=blue!8] (self) at (-3.30,0.55) {自身变换\\$W_{\rm self}\mathbf x_i$};
\node[op,fill=orange!10] (neigh) at (-3.30,-0.55) {邻居变换\\$W_{\rm neigh}\sum_j\mathbf x_j$};
\node[merge,fill=purple!9] (add) at (-0.50,0) {逐元素相加\\$\sigma(\mathbf a+\mathbf b)$};
\node[read] (readout) at (1.80,0) {共享\\读出层};
\draw[arrow] (xi) -- (self);
\draw[arrow] (xj) -- (neigh);
\draw[arrow] (self) -- (add);
\draw[arrow] (neigh) -- (add);
\draw[arrow] (add) -- (readout);
\node[align=left,text width=50mm,font=\footnotesize] at (5.20,0)
  {自身项与邻居项分别变换，但在中间层相加；高度数节点可能引入量级失衡与表示干扰。};

% Panel (b): two-channel representation.
\begin{scope}[yshift=-3.30cm]
\node[font=\bfseries,anchor=west] at (-6.20,1.45) {(b) 自身--邻居双通道};
\node[input] (xi2) at (-5.60,0.55) {$\mathbf x_i$};
\node[input] (xj2) at (-5.60,-0.55) {$\{\mathbf x_j\}$};
\node[op,fill=blue!8] (self2) at (-3.30,0.55) {自身通道\\$f_{\rm self}(\mathbf x_i)$};
\node[op,fill=orange!10] (neigh2) at (-3.30,-0.55) {邻居通道\\$\sum_j f_{\rm neigh}(\mathbf x_j)$};
\node[merge,fill=green!9] (cat) at (-0.50,0) {表示拼接\\$[\mathbf h_i^{\rm self};\mathbf h_i^{\rm neigh}]$};
\node[read] (readout2) at (1.80,0) {非线性\\读出层};
\draw[arrow] (xi2) -- (self2);
\draw[arrow] (xj2) -- (neigh2);
\draw[arrow] (self2) -- (cat);
\draw[arrow] (neigh2) -- (cat);
\draw[arrow] (cat) -- (readout2);
\node[align=left,text width=50mm,font=\footnotesize] at (5.20,0)
  {两类信息在读出前保持分离；邻居通道采用求和，以保留局部交互的累积强度。};
\end{scope}
\end{tikzpicture}
